\section{Überblick über die erhobenen Daten}
\label{sec:05-01_data-overview}

Die empirische Untersuchung stützt sich auf eine kombinierte Datenerhebung aus Usability-Tests, standardisierten Skalenfragen sowie leitfadengestützten Interviews mit offenen und halb-offenen Elementen. Ziel der Erhebung war es, objektive Nutzungsdaten mit subjektiven Einschätzungen der Teilnehmenden zu verbinden. Sämtliche Rohdaten, tabellarischen Auswertungen und Interviewprotokolle sind im Anhang dokumentiert und den jeweiligen Interview-IDs (TN-01 bis TN-10) eindeutig zugeordnet (vgl.~Anhang).

An der Untersuchung nahmen insgesamt zehn Personen (\(n = 10\)) teil. Das Alter der Teilnehmenden lag zwischen 22 und 52 Jahren, mit einem arithmetischen Mittelwert von 35{,}2 Jahren und einem Median von 36{,}5 Jahren. Sechs Teilnehmende waren männlich, vier weiblich. Sieben Personen gehörten der betrachteten Produktorganisation an, drei Personen waren extern. Innerhalb der organisationalen Stichprobe waren unterschiedliche Rollen vertreten, darunter sowohl operative Tätigkeiten als auch koordinierende und leitende Funktionen. Eine detaillierte Übersicht zu soziodemografischen Merkmalen, Rollenprofilen und Organisationszugehörigkeit ist im Anhang dargestellt (vgl.~Anhang).

Die Nutzungshäufigkeit des bestehenden Wikis zeigte eine heterogene Verteilung. Fünf Teilnehmende gaben an, das bisherige Wiki mindestens einmal pro Woche zu nutzen, drei Personen nutzten es nur gelegentlich (monatlich), während zwei Personen angaben, es bislang nicht genutzt zu haben. Externe Teilnehmende gehörten ausschließlich zu den sporadischen Nutzenden oder Nicht-Nutzenden. Ein klarer Zusammenhang zwischen Rolle oder hierarchischer Position und Nutzungshäufigkeit ließ sich in der Stichprobe nicht erkennen (vgl.~Anhang).

Die digitale Selbsteinschätzung der Teilnehmenden fiel durchgängig hoch aus. Alle zehn Personen stimmten der Aussage zu, sich im Umgang mit digitalen Tools sicher zu fühlen. Sieben Teilnehmende wählten eine überwiegend zustimmende Antwortoption, drei eine eher zustimmende; neutrale oder ablehnende Antworten traten nicht auf. Damit weist die Stichprobe ein insgesamt hohes digitales Kompetenzniveau auf, wodurch Unterschiede in der Aufgabenbearbeitung oder in der subjektiven Bewertung nicht auf mangelnde digitale Erfahrung zurückgeführt werden können (vgl.~Anhang).

Im Rahmen der Usability-Tests bearbeitete jede Person fünf definierte Aufgaben sowohl im bestehenden als auch im neu gestalteten Wiki. Insgesamt liegen damit 100 Einzelmessungen vor. Erfasst wurden die Bearbeitungszeit in Sekunden, die Anzahl der Fehler sowie die Anzahl der Klicks pro Aufgabe. Über alle Aufgaben hinweg reduzierte sich die durchschnittliche Bearbeitungszeit von 60{,}68 Sekunden in der alten Version auf 19{,}18 Sekunden in der neuen Version; der Median sank von 48 Sekunden auf 17 Sekunden. Parallel dazu verringerte sich die Gesamtzahl der Fehler von 100 auf 15. Diese Effekte traten bei allen Teilnehmenden auf, unabhängig von Nutzungshäufigkeit oder Rolle (vgl.~Anhang).

Ergänzend zu den quantitativen Nutzungsdaten wurden leitfadengestützte Interviews durchgeführt. Der Leitfaden umfasste vier thematische Blöcke zu wahrgenommener Nützlichkeit und Gebrauchstauglichkeit, Rollen- und Persona-Passung, Akzeptanz und Nutzungsintention sowie offenem Feedback. Die geschlossenen Fragen lieferten skalierte Bewertungen, während halb-offene und offene Fragen zur Begründung, Vertiefung und Ergänzung genutzt wurden. Die qualitativen Daten wurden fallbezogen dokumentiert und anschließend mittels qualitativer Inhaltsanalyse nach Mayring ausgewertet. Die deduktiv-induktiv entwickelte Kategorienstruktur sowie die vollständigen Kodierungen sind im Anhang nachvollziehbar dargestellt (vgl.~Anhang).

Insgesamt liegt damit eine konsistente und umfangreiche Datenbasis vor, die objektive Leistungskennzahlen, subjektive Bewertungen und qualitative Aussagen systematisch miteinander verbindet. Diese Daten bilden die Grundlage für die folgenden Unterkapitel, in denen die Ergebnisse entlang der Forschungsfragen zu Informationsarchitektur, Rollenpassung, UX-Strategien und Akzeptanzfaktoren detailliert dargestellt werden.